\section{Universal multivariate Newton atlas}\label{sec:mv1}

Return to the actual coordinates \(d_1,\ldots,d_4,t\).  By
\cref{cor:sixsupport}, every non-zero moment monomial \(d^\gamma t^k\) with
\(\gamma\ne0\) satisfies \(k\ge|\gamma|+2\), and the boundary consists only of
pure powers.

\begin{theorem}[Universal Newton polyhedron]\label{thm:universalnewton}
The Newton polyhedron of \(K(d,t)\) at the dark origin is
\begin{equation}\label{eq:universalNP}
\conv\{(e_1,3),(e_2,3),(e_3,3),(e_4,3)\}+\R_{\ge0}^{5}.
\end{equation}
At the four vertices the raw moment coefficients are \(+1,-1,+1,-1\), and the
amplitude coefficients are \(i/6,-i/6,i/6,-i/6\).
\end{theorem}
\begin{proof}
For \(\gamma=0\), all coefficients vanish by the all-depth baseline-darkness
theorem, so no unperturbed exponent enters the universal support.  The
\(j=1\) case of \cref{eq:minlayer} supplies the four non-zero vertices.  For
any other exponent \((\gamma,k)\) with \(\gamma\ne0\), choose \(r\) with
\(\gamma_r>0\).  Then \(k\ge|\gamma|+2\ge3\), and
\((e_r,3)\le(\gamma,k)\) coordinatewise.  Hence all actual support lies in the
upper closure of the four vertices.  The amplitude coefficients follow from
\((-i)^3/3!=i/6\).
\end{proof}

For a positive weight \(\omega=(\omega_1,\ldots,\omega_4,\tau)\), the exposed
vertices are indexed by the non-empty subset
\[
I(\omega)=\{r:\omega_r=\min_j\omega_j\}.
\]
Every non-empty subset occurs, giving exactly \(2^4-1=15\) relatively open
cells.  On a cell \(I\), the face polynomial is
\begin{equation}\label{eq:mvface}
\frac{i}{6}t^3\sum_{r\in I}\sigma_r d_r,\qquad
(\sigma_1,\sigma_2,\sigma_3,\sigma_4)=(1,-1,1,-1).
\end{equation}
Over the complex torus, every face with at least two vertices has a cancellation
hypersurface.  Over the positive real torus, a face cancels only when \(I\)
contains both signs.

\begin{figure}[ht]
\centering\begin{tikzpicture}[scale=1.05,>=Latex]
\coordinate (A) at (0,2.8); \coordinate (B) at (-2,-1.1); \coordinate (C) at (2,-1.1); \coordinate (D) at (0,.1);
\fill[blue!7] (A)--(B)--(C)--cycle;
\draw[thick] (A)--(B)--(C)--cycle; \draw[thick] (A)--(D)--(B); \draw[thick] (D)--(C);
\draw[dashed] (B)--(C); \draw[thick] (D)--(A);
\fill[blue!70!black] (A) circle (2.5pt); \node[above,font=\small] at (A) {$(e_1,3)$};
\fill[blue!70!black] (B) circle (2.5pt); \node[below left,font=\small] at (B) {$(e_2,3)$};
\fill[blue!70!black] (C) circle (2.5pt); \node[below right,font=\small] at (C) {$(e_3,3)$};
\fill[blue!70!black] (D) circle (2.5pt); \node[right,font=\small] at (D) {$(e_4,3)$};
\node[align=center,font=\small] at (4.3,.8) {$15$ non-empty subsets\\$4+6+4+1$};
\draw[->,thick] (3.1,.8)--(2.1,.7);
\end{tikzpicture}

\caption{A tetrahedral projection of the four Newton vertices.  Its non-empty
vertex subsets encode the 15 positive normal cells; this is a combinatorial
projection, not a Euclidean embedding of the five-dimensional polyhedron.}
\label{fig:mv1cells}
\end{figure}

\subsection{Moment--amplitude face mismatch}

The moment and amplitude have the same exponent support, because the multiplier
\((-i)^k/k!\) is never zero.  Their face cancellation equations need not agree
when a pullback makes distinct time depths tie.  A reachable two-term face with
parent points \((e_1,4)\) and \((e_2,3)\) has raw coefficients \(2,-1\).  Its
moment and amplitude face forms are
\[
W_M=2c_1-c_2,\qquad W_A=\frac{c_1}{12}-\frac{i c_2}{6}.
\]
Thus equality of supports does not imply equality of cancellation loci.

\subsection{Fixed positive rational weights}

The following corrects an overly pessimistic interpretation of the finite
universal atlas obstruction.

\begin{theorem}[Fixed-weight coefficient-stratified termination]\label{thm:noetherian}
For every fixed positive rational weight, after finite ramification where
needed, the pulled amplitude has a finitely generated coefficient ideal, so a
finite coefficient-stratified exact-zero/first-survivor decision exists.  No
effective uniform depth bound or finite global recursive atlas over all
positive weights has been proved.
\end{theorem}
\begin{proof}
Clear denominators in the fixed rational weight by a ramification \(s=u^m\).
After substituting
\(d_r=c_ru^{b_r}\) and \(t=c_tu^a\), write the pulled amplitude as
\[
K_\omega(u;c)=\sum_{n\ge0}A_n(c)u^n,
\]
where each \(A_n\) is a polynomial in the leading constants over the base field
(or a Laurent polynomial if constants are restricted to the torus).  The ring
\(\mathbb K[c_1,\ldots,c_4,c_t]\), and likewise its Laurent localisation, is
Noetherian \parencite{AtiyahMacdonald1969}.  Therefore the ideal
\(J=(A_n:n\ge0)\) is generated by a finite subset, say with indices at most
\(N\).  If \(A_0(c)=\cdots=A_N(c)=0\), then every \(A_n(c)=0\), so the pullback
is exactly zero.  Otherwise its first survivor occurs among those finitely many
indices.  This proves existence for the fixed weight.  The argument neither
computes a useful \(N\) nor bounds \(N\) uniformly as the weight varies.
\end{proof}

The theorem concerns a coefficient-stratified decision, not a weight-only
decision.  Delayed cancellations can make the necessary depth arbitrarily
large across families of weights or germs; that fact is compatible with
Noetherianity at each fixed weight.
